\section{引言}

卫星海洋遥感大数据与自主水下航行器的结合，正在将海洋环境监测从按预设航线走航转变为由先验地图引导的现场自适应巡测\cite{fossum2019}。在这一观测范式中，多年累积的卫星遥感产品提供了大尺度的海洋环境背景与演变统计，而自主水下航行器则深入云层遮挡与遥感平滑无法企及的现场水体采集高分辨率样本。然而机动观测平台受制于有限的机载电池能量、嵌入式计算速率与固定的作业时间窗口\cite{xiao2020thinking, magnino2026metareasoning}。航行器在现场获取每一条新观测后，既可以调用机载高功耗芯片开展密集的未来航线推演以优化后续探索路径，也可以维持低功耗处理器后台规划而将宝贵的电能与时间用于长距离巡航采样或原位重复测量\cite{horvitz1987reasoning, russell1991right, shen2025thinking}。由于内部航线计算与外部物理观测共同消耗同一块电池和同一个任务窗口，何时应当继续计算、何时应当继续观测，构成了大数据驱动自主海洋观测系统的核心控制问题\cite{hunt2021formalizing}。

东海有害赤潮应急巡测集中体现了先验卫星底图失真、水下通信受阻与机载资源受限带来的决策困难\cite{fossum2019, lou2014diurnal}。根据2023年中国海洋灾害公报，我国海域全年发现赤潮46次，累计面积约$1500\text{ km}^2$，其中有毒有害赤潮达29次\cite{mnr2024bulletin}。在浙江近海台州、温州与宁波外海，晚春至初夏常暴发高密度的东海原甲藻赤潮\cite{zhao2025distinguishing}，且表层藻华斑块在潮汐作用下发生剧烈的往复漂移与形态演变\cite{lou2014diurnal}。赤潮暴发时，应急指挥中心通常向$80\text{ km}\times 80\text{ km}$的目标海域派出一台百米级自主水下航行器开展24小时紧急巡测，要求航行器尽早定位叶绿素浓度最高的核心水体，并随时响应指挥中心对最浓处位置的查询\cite{remus100}。受多星遥感产品合成延迟、云区经验正交函数插补平滑\cite{beckers2003dineof}、水色反演固有偏差\cite{ye2021cocts}以及调查船航渡时间的共同影响，航行器下水前装入的往往是3天前的旧卫星图。加之航行器在水下作业期间几乎无法与母船或岸站实时通信，全部航线重规划与资源调配均须由机载计算机在航行途中独立完成\cite{jetsonorinnx}。

传统的海洋巡测规则与预设时间分配策略无法在缺乏历史任务数据的条件下适配不同的机载硬件。一方面，海洋调查中广泛采用的割草机式全覆盖\cite{galceran2013coverage}、贪心追踪以及基于高斯过程的上置信界采样\cite{srinivas2012gpucb}和汤普森采样\cite{russo2018thompson}等启发式自适应规则，均将每步决策的计算开销视为零或固定常数，忽略了多步前瞻航线规划对规避局部次优峰值的关键作用\cite{hollinger2014sampling}。另一方面，若在多步规划器之上采用固定的时间分配策略，例如全程仅靠低功耗处理器后台规划的“一直走”或周期性开启高功耗芯片的固定比例规则，其表现高度依赖于机载处理器的推演速度与芯片功耗。当航行器更换嵌入式计算模块或电池配置时，唯有依靠该型硬件在大量历史赤潮任务中的实测表现才能选出合适的固定策略，而在新平台首次投入应急巡测时并不具备这一历史任务档案。

从历史巡测任务中直接训练策略选择器同样受困于高昂的数据采集成本与剧烈的任务间环境波动。近年来机器人元推理研究尝试利用历史任务数据训练回归或分类模型，根据出发前的环境特征为每次任务自动选择计算预算或时间分配策略\cite{sung2021quit, magnino2026metareasoning}。然而海洋赤潮现场的初值格局、峰值位置与航行器入水起点千差万别，不同任务窗口之间的天然遗憾波动远远大于几种优良用时策略之间的平均性能差距。因此，基于旧卫星图静态特征训练的选择器极易被任务间的背景噪声淹没，不仅需要上百次历史巡测任务作为训练样本，而且在未知测试年份上的泛化表现依然无法超越简单的固定策略。

人工智能领域的经典在线单步元推理虽无需历史任务训练，但因轻信规划器在搜索样本上的自评得分而不可避免地陷入过度规划\cite{horvitz1987reasoning, russell1991right, cuadron2025overthinking}。经典元推理在每个控制周期通过比较“多算一步带来的航线得分增量”与“耗电及延误造成的代价”来在线决定是否开启高功耗芯片或停下复测\cite{lin2015metareasoning, elboher2023formal}。由于航行器无法预知真实海洋，规划器只能在一组由旧卫星图生成的想象海洋上搜索航线。当规划器在同一批有限的想象海洋上既搜索最优航路点又评估该航线的预期收益时，优化过程会精细迎合这批特定样本中的随机正涨落，产生决策分析中的“优化器诅咒”\cite{smith2006curse, vanhasselt2010double}。即使额外推演在真实海洋中早已不再改善航线质量，规划器看到的内部自评得分仍随推演次数增加而持续虚高，进而误导单步元推理频繁开启高功耗芯片并中途停下漂流\cite{jiang2015horizon, cuadron2025overthinking}。

针对上述症结，本文的核心思想是：航行器绝不能依据规划器自身的搜索得分来决定是否继续计算，而应利用多年卫星海洋大数据标定误差的留出想象海洋作为独立标尺，在未参与航线搜索的想象海洋上评判计算与观测的真实结果。具体而言，本文利用八年东海卫星叶绿素数据标定滞后底图的误差统计，围绕出发前旧图生成48个随现场荧光计采样经集合调整卡尔曼滤波\cite{anderson2001eakf, gaspari1999correlation}同步更新的想象海洋，并将其严格划分为32个搜索用海洋与16个留出海洋。基于这一“搜索与评判分离”准则，本文提出了两种无需历史巡测任务的控制方式：下水前，利用母船或岸基计算机以4个留出海洋为假想真值，按机载算力与功耗将五种候选策略各完整预演一遍24小时任务并择优执行；航行中，若需按周期在线决策，则在单步元推理中将计算与采样的价值改由16个留出海洋评估，从而在离线与在线两个层面根除自评乐观偏差。

本文的主要贡献归纳如下：
\begin{itemize}
  \item \textbf{提出独立于搜索过程的计算价值留出评判准则：}揭示了经典元推理在未知海洋巡测中过度计算的根源在于规划器对搜索用海洋的过拟合，提出将多年卫星遥感数据标定的想象海洋划分为搜索用海洋与留出海洋，以搜索未曾接触的留出海洋上的结果作为衡量计算收益的无偏标尺。
  \item \textbf{构建面向硬件自适应的出发前全任务预演与在线留出估值方法：}设计了结合局地化集合同化与交叉熵航线规划的五种用时策略，提出在母船或岸基计算机上以留出海洋为假想真值开展全任务出发前预演，并将同一留出评判机制嵌入在线单步元推理，无需任何历史巡测任务即可适配不同机载处理器速度与功耗。
  \item \textbf{提供基于八年东海赤潮窗口与双真值模型的系统实证及边界刻画：}在3197天卫星数据构建的422个赤潮窗口、四种机载硬件与两种真值模型上开展了配对重抽样检验，证实深度时空模型无法替代现场采样、留出评判可将跨三个半数量级推演预算的价值估计偏差控制在0.005以内、停下复测在全参数空间均得不偿失，且出发前预演与在线留出估值均达到依赖历史任务方法的同等水平。
\end{itemize}

本文其余章节安排如下：第2节回顾元推理、海洋自适应采样及优化器诅咒与集合数据同化的相关研究；第3节建立东海赤潮应急巡测的场景模型、航行器机电功耗模型与时间平均遗憾目标，并剖析问题的三重内在矛盾；第4节阐述数据标定想象海洋、搜索与评判分离、出发前预演及在线留出估值方法；第5节报告遥感数据统计、自评偏差验证、主实验对比、作用机制、消融实验与适用边界相图，并讨论研究局限性；第6节总结全文结论。
